\section{Immersions of schemes}
\label{section-immersions}

\noindent
In Lemma \ref{lemma-open-subspace-scheme} we saw that any open subspace
of a scheme is a scheme. Below we will prove that the same holds for
a closed subspace of a scheme.

\medskip\noindent
Note that the notion of a quasi-coherent sheaf of $\mathcal{O}_X$-modules
is defined for any ringed space $X$ in particular when $X$ is a scheme.
By our efforts in Section \ref{section-quasi-coherent-affine} we
know that such a sheaf is on any affine open $U \subset X$
of the form $\widetilde M$ for some $\mathcal{O}_X(U)$-module $M$.

\begin{lemma}
\label{lemma-closed-subspace-scheme}
Let $X$ be a scheme. Let $i : Z \to X$ be a closed immersion
of locally ringed spaces.
\begin{enumerate}
\item The locally ringed space $Z$ is a scheme,
\item the kernel $\mathcal{I}$ of the map
$\mathcal{O}_X \to i_*\mathcal{O}_Z$ is a quasi-coherent
sheaf of ideals,
\item for any affine open $U = \Spec(R)$ of $X$
the morphism $i^{-1}(U) \to U$ can be identified with
$\Spec(R/I) \to \Spec(R)$ for some ideal $I \subset R$, and
\item we have $\mathcal{I}|_U = \widetilde I$.
\end{enumerate}
In particular, any sheaf of ideals locally generated by sections
is a quasi-coherent sheaf of ideals (and vice versa),
and any closed subspace of $X$ is a scheme.
\end{lemma}

\begin{proof}
Let $i : Z \to X$ be a closed immersion.
Let $z \in Z$ be a point. Choose any affine open
neighbourhood $i(z) \in U \subset X$. Say $U = \Spec(R)$.
By Lemma \ref{lemma-closed-immersion-affine-case} we know
that $i^{-1}(U) \to U$ can be identified with the morphism
of affine schemes $\Spec(R/I) \to \Spec(R)$.
First of all this implies that $z \in i^{-1}(U) \subset Z$ is an
affine neighbourhood of $z$. Thus $Z$ is a scheme. Second
this implies that $\mathcal{I}|_U$ is $\widetilde I$.
In other words for every point $x \in i(Z)$ there exists an
open neighbourhood such that $\mathcal{I}$ is quasi-coherent in
that neighbourhood. Note that $\mathcal{I}|_{X \setminus i(Z)}
\cong \mathcal{O}_{X \setminus i(Z)}$. Thus the restriction
of the sheaf of ideals is quasi-coherent on $X \setminus i(Z)$
also. We conclude that $\mathcal{I}$ is quasi-coherent.
\end{proof}

\begin{definition}
\label{definition-immersion}
Let $X$ be a scheme.
\begin{enumerate}
\item A morphism of schemes is called an {\it open immersion}
if it is an open immersion of locally ringed spaces (see
Definition \ref{definition-immersion-locally-ringed-spaces}).
\item An {\it open subscheme} of $X$ is an open subspace of $X$
in the sense of Definition \ref{definition-open-subspace}; an open subscheme
of $X$ is a scheme by Lemma \ref{lemma-open-subspace-scheme}.
\item A morphism of schemes is called a {\it closed immersion}
if it is a closed immersion of locally ringed spaces (see
Definition \ref{definition-closed-immersion-locally-ringed-spaces}).
\item A {\it closed subscheme} of $X$ is a closed subspace of $X$
in the sense of Definition \ref{definition-closed-subspace}; a closed subscheme
is a scheme by Lemma \ref{lemma-closed-subspace-scheme}.
\item A morphism of schemes $f : X \to Y$ is called an {\it immersion},
or a {\it locally closed immersion} if it can be factored as
$j \circ i$ where $i$ is a closed immersion and $j$ is an open
immersion.
\end{enumerate}
\end{definition}

\noindent
It follows from the lemmas in Sections \ref{section-open-immersion} and
\ref{section-closed-immersion} that any open (resp.\ closed) immersion of
schemes is isomorphic to the inclusion of an open (resp.\ closed) subscheme
of the target.

\medskip\noindent
Our definition of a closed immersion is halfway between Hartshorne
and EGA. Hartshorne defines a closed immersion as a morphism $f : X \to Y$ of
schemes which induces a homeomorphism of $X$ onto a closed subset of
$Y$ such that $f^\# : \mathcal{O}_Y \to f_*\mathcal{O}_X$ is surjective, see
\cite[Page 85]{H}. We will show this is equivalent to our notion in
Lemma \ref{lemma-characterize-closed-immersions}. In \cite{EGA},
Grothendieck and Dieudonn\'e first define closed subschemes
via the construction of Example \ref{example-closed-subspace}
using quasi-coherent sheaves of ideals and then define a closed immersion
as a morphism $f : X \to Y$ which induces an isomorphism with a
closed subscheme. It follows from Lemma \ref{lemma-closed-subspace-scheme}
that this agrees with our notion.

\medskip\noindent
Pedagogically speaking the definition above is a disaster/nightmare.
In teaching this material to students, we have found it often
convenient to define a closed immersion as an affine morphism
$f : X \to Y$ of schemes such that $f^\# : \mathcal{O}_Y \to f_*\mathcal{O}_X$
is surjective. Namely, it turns out that the notion of an affine
morphism (Morphisms, Section \ref{morphisms-section-affine})
is quite natural and easy to understand.

\medskip\noindent
For more information on closed immersions we suggest the reader visit
Morphisms, Sections \ref{morphisms-section-closed-immersions} and
\ref{morphisms-section-closed-immersions-quasi-coherent}.

\medskip\noindent
We will discuss locally closed subschemes and immersions
at the end of this section.

\begin{remark}
\label{remark-not-reverse-open-closed}
If $f : X \to Y$ is an immersion of schemes, then it is in general
not possible to factor $f$ as an open immersion followed
by a closed immersion. See Morphisms, Example \ref{morphisms-example-thibaut}.
\end{remark}

\begin{lemma}
\label{lemma-immersion-when-closed}
Let $f : Y \to X$ be an immersion of schemes. Then $f$ is a closed
immersion if and only if $f(Y) \subset X$ is a closed subset.
\end{lemma}

\begin{proof}
If $f$ is a closed immersion then $f(Y)$ is closed by definition.
Conversely, suppose that $f(Y)$ is closed. By definition
there exists an open subscheme $U \subset X$ such that $f$ is the composition
of a closed immersion $i : Y \to U$ and the open immersion
$j : U \to X$. Let $\mathcal{I} \subset \mathcal{O}_U$ be the
quasi-coherent sheaf of ideals associated to the closed immersion
$i$. Note that
$\mathcal{I}|_{U \setminus i(Y)}
= \mathcal{O}_{U \setminus i(Y)}
= \mathcal{O}_{X \setminus i(Y)}|_{U \setminus i(Y)}$.
Thus we may glue (see Sheaves, Section \ref{sheaves-section-glueing-sheaves})
$\mathcal{I}$ and $\mathcal{O}_{X \setminus i(Y)}$ to a sheaf of
ideals $\mathcal{J} \subset \mathcal{O}_X$. Since every
point of $X$ has a neighbourhood where $\mathcal{J}$ is
quasi-coherent, we see that $\mathcal{J}$ is quasi-coherent
(in particular locally generated by sections).
By construction $\mathcal{O}_X/\mathcal{J}$ is supported
on $U$ and, restricted there, equal to $\mathcal{O}_U/\mathcal{I}$.
Thus we see that the closed subspaces associated to $\mathcal{I}$
and $\mathcal{J}$ are canonically isomorphic, see
Example \ref{example-closed-subspace}.
In particular the closed subspace of $U$ associated to $\mathcal{I}$
is isomorphic to a closed subspace of $X$.
Since $Y \to U$ is identified with the closed subspace
associated to $\mathcal{I}$, see Lemma \ref{lemma-closed-immersion},
we conclude that $Y \to U \to X$
is a closed immersion.
\end{proof}

\noindent
Let $f : Y \to X$ be an immersion.
Let $Z = \overline{f(Y)} \setminus f(Y)$ which is a closed subset of $X$.
Let $U = X \setminus Z$.
The lemma implies that $U$ is the biggest open subspace of $X$
such that $f : Y \to X$ factors through a closed immersion into $U$.
We define a {\it locally closed subscheme of $X$} as a pair
$(Z, U)$ consisting of a closed subscheme $Z$ of an open subscheme $U$
of $X$ such that in addition $\overline{Z} \cup U = X$. We usually
just say ``let $Z$ be a locally closed subscheme of $X$'' since we may
recover $U$ from the morphism $Z \to X$. The above
then shows that any immersion $f : Y \to X$ factors uniquely
as $Y \to Z \to X$ where $Z$ is a locally closed subspace of $X$
and $Y \to Z$ is an isomorphism.

\medskip\noindent
The interest of this is that the collection of locally closed subschemes
of $X$ forms a set. We may define a partial ordering on this set,
which we call inclusion for obvious reasons. To be explicit, if
$Z \to X$ and $Z' \to X$ are two locally closed subschemes of $X$, then
we say that {\it $Z$ is contained in $Z'$} simply if the morphism
$Z \to X$ factors through $Z'$. If it does, then of course $Z$ is identified
with a unique locally closed subscheme of $Z'$, and so on.

\medskip\noindent
For more information on immersions, we refer the reader to
Morphisms, Section \ref{morphisms-section-immersions}.










